% Authored proof supplement, not a generated surface. Included by contracts.
\subsection*{Round 4: proved subgates and binding residuals (2026-09-05)}
\label{sec:toe-round4-proofs}
The following results concern the explicitly stated regulators and algebras.
They do not identify a single complete physical $3{+}1$D parent. All T1--T8
and \texttt{TFPT.TOE.COMPLETE.01} remain \mO. The finite regressions are
checks of the proofs below, not substitutes for their quantifiers.

\paragraph{T2: uniform relaxed TEL-B theorem.}
\texttt{SEAM.MMST.TELB.NORM.01}.
For the frozen $M_{\rm mass}=1$, $N_y=8$ collar, the original remainder
obeys
\[
 \boxed{\|R_N\|_{\rm HS}<2.995906<3\quad(N\ge16\text{ even}).}
\]
The interval inputs are v1022 (TV and potential), v1025 (C1/C2a), and
\veri{v1026\_\allowbreak telb\_\allowbreak round4\_\allowbreak closure.py}
(CF, DG, signed CROSS, validated Fourier core and all-volume aliases).
This is not a proof of the historical sharper $2.11$ or Bound-B target
$0.62$. In particular a limiting Fourier-core bound is not a finite-grid
remainder bound.

\paragraph{CF: cell covariance and the periodic corner.}
Put $h=N^{-1}$, $A=(2\pi^2)^{-1}$, $J=N/2$,
$I_N=\{-J,\ldots,J-1\}$ and let bars denote normalized means over
$[kh,(k+1)h]$. Define
\[
 \begin{split}
 f(u)&=-\log|u|-\pi\mathbf1_{u<0},\qquad
 \Phi(u)=\log\frac{|u|}{2|\sin\pi u|}-\pi u,\qquad F=f+\Phi,\\
 c(d)&=d^{-1}-\pi\cot\pi d,\quad c(0)=0,\qquad
 Q_g(u,v)=\frac{g(v)-g(u)}{v-u}.
 \end{split}
\]
Partial fractions give $\Phi'=c-\pi$, $L=\sup|\Phi'|=\pi+2$,
$m=\inf\Phi''=\pi^2/3$, $M=\sup\Phi''=\pi^2-4$ and
$D=M-m=2\pi^2/3-4$, with $\int|\Phi'''|=2D$.
The identities used are the standard cotangent and polygamma identities
\href{https://dlmf.nist.gov/4.22}{DLMF 4.22} and
\href{https://dlmf.nist.gov/5.15}{DLMF 5.15}; no asymptotic truncation is needed.
The real kernel cancellation is exact:
\[
 r(u,v)=A\{c(v-u)(F(v)-F(u))-Q_\Phi(u,v)
                  -i\pi(F(v)-F(u))\}.
\]
For independent uniform cell points $U,V$, set $d_0=(l-k)h$ and
$\delta=V-U-d_0$. Then
\[
 (\mathcal K_N)_{kl}=h\mathbb E r(U,V),\quad
 (\mathcal K'_N)_{kl}=hA\left[c(d_0)(\bar F_l-\bar F_k)
 -\frac{\bar\Phi_l-\bar\Phi_k}{d_0}-i\pi(\bar F_l-\bar F_k)\right].
\]
Thus the imaginary part freezes exactly. If $v_k=\operatorname{Var}F(U_k)$,
$|a'|\le a_1$, $|a''|\le a_2$, the independent-copy variance identity,
$\mathbb E\delta=0$, and $\mathbb E\delta^2=h^2/6$ give
\[
 |\mathbb E[(a(d_0)-a(V-U))(F(V)-F(U))]|
 \le\frac{a_1h}{\sqrt6}\sqrt{v_k+v_l}
       +\frac{a_2h^2}{12}|\bar F_l-\bar F_k|.                 \tag{R4.1}
\]
For $|g''|\le B$, the two coordinate Lipschitz constants of $Q_g$ are
at most $B/2$, so $\operatorname{Var}Q_g\le B^2h^2/24$ and
$|\mathbb E(\delta Q_g)|\le Bh^2/12$.
The two singular log cells have variance one; all remaining log-cell
variances are at most $1/(12n^2)$. Therefore
\[
 (\textstyle\sum_k v_k)^{1/2}\le
 V(h):=\sqrt{2+\zeta(2)/6}+L\sqrt{h/12},\quad
 \textstyle\sum_{k,l}|\bar F_l-\bar F_k|^2\le2N^2\zeta(2).
\]
The last bound follows from $\operatorname{Var}F=\pi^2/6$ and contraction
of cell conditional expectation. Split the nonzero lags into central
($|l-k|\le N/2$), middle ($N/2<|l-k|\le3N/4$), and far regions.
For the first two regions (R4.1) gives
\[
 \begin{split}
 C_{\rm central}&\le\frac{Ac'(9/16)hV(h)}{\sqrt3}
  +\frac{Ac''(9/16)h^{3/2}\sqrt{2\zeta(2)}}{12}<0.027267,\\
 C_{\rm middle}&\le\frac{Ac'(13/16)h^{3/2}}{\sqrt6}
       \left(\sqrt{H_{J-1}/3}+L/\sqrt{24}\right)
  +\frac{Ac''(13/16)h^{3/2}\sqrt{2\zeta(2)}}{12}<0.055905.
 \end{split}
\]
In the middle count, negative index $-n$ has $n-1$ partners and positive
index $s$ has $s$ partners; the singular cells have none. Consequently
the log-variance sum is at most $H_{J-1}/3$ and the smooth part at most
$L^2/24$. The sequence $H_{J-1}/J^3$ decreases for $J\ge8$, by its
one-step recurrence, so the displayed envelopes are evaluated at $N=16$.

For the far seam let $r=N-(l-k)$, $t=1-(V-U)$ and $t_0=rh$.
The short periodic continuation satisfies $F(V)-F(U)=-tQ$ with
$|F'|\le L_F=2\pi$ and $|F''|\le M_F=2\pi^2$. In
$c(d)=(1-d)^{-1}+a(d)$ the apparent pole cancels exactly:
\[
 \mathbb E[(t_0^{-1}-t^{-1})(F(V)-F(U))]
       =-t_0^{-1}\mathbb E[(t-t_0)Q].
\]
There are $r$ pairs per orientation. With
$a_1=(16/11)^2+c'(5/16)$, $a_2=2(16/11)^3+c''(5/16)$,
\[
 C_{\rm far}\le AM_Fh^{3/2}\sqrt{H_{\lfloor N/4\rfloor}/72}
 +Aa_1L_Fh^{3/2}/6+Aa_2L_F(5/16)h^{3/2}/12<0.009280.
\]
Here $H_{\lfloor N/4\rfloor}/N^3\le H_4/16^3$ follows by grouping four
successive sizes. The remaining smooth quotient contributes
$C_\Phi\le AMh\sqrt{\zeta(2)/72}<0.002810$. The three lag regions are
HS-orthogonal; only the last contribution uses the triangle inequality:
\[
 \sqrt N\| (\mathcal K'_N-\mathcal K_N)_{\rm off}\|_{\rm HS}
 \le\sqrt{C_{\rm central}^2+C_{\rm middle}^2+C_{\rm far}^2}+C_\Phi
 <0.065697.                                                   \tag{R4.2}
\]

\paragraph{DG: original diagonal, not a replacement recurrence.}
The csc-squared finite product identity and trigamma recurrence for the
original remainder give, for $R_k=(\mathrm{Res}_N)_{kk}$,
\[
 R_{k+1}-R_k=-Ah^2\Phi''((k+3/4)h),\qquad
 \sum_kR_k=A\pi+A\tau(J),
\]
where $T(x)=(x-1)\psi_1(x)+\psi(x)$ and
$\tau(J)=T(J+5/4)-T(J+3/4)$. Since
$T'(x)=2\sum_{n\ge0}(n+1)/(n+x)^3\le4/J$ for $x\ge J\ge8$,
$0\le\tau(J)\le2/J=4h$. The $3/4$-point Peano kernel and
$h\sum_k\Phi'(kh)=-\pi-2h$ imply
\[
 |NR_k+A\Phi'(kh)|\le Ah(2+3D/2).                           \tag{R4.3}
\]
On each diagonal cell the imaginary mean vanishes,
$|Q_\Phi-\Phi'(kh)|\le Mh$, and
$|c(d)|\le C_s|d|$, $C_s=(\pi^2/3)/(1-h^2)$.
The worst logarithmic integral occurs adjacent to zero and is exactly
$\mathbb E[|V-U||\log|V/U||]\le h/2$; also
$\mathbb E|V-U|^2=h^2/6$. Thus
\[
 \sqrt N\|(\mathrm{Res}_N-\mathcal K_N)_{\rm diag}\|_{\rm HS}
 \le Ah(2+3D/2)+AMh+AC_s(h/2+Lh^2/6)<0.042959.              \tag{R4.4}
\]
All scalar envelopes increase with $h\le1/16$. With the independently
certified potential contribution $0.399224424139108$, orthogonality gives
$c_E<0.466902$, below the required $0.542702$.

\paragraph{CROSS: the sufficient one-sided theorem.}
Let $W_\infty$ be the unitary fractional shift with entries
$e^{-i\pi/4}\sin(\pi/4)/[\pi(k-m-1/4)]$, let $P$ restrict to $I_N$,
$w=P W_\infty e_0$, $B=W_\infty P_{(-\infty,-1]}W_\infty^*$,
$t_N=\|(1-P)W_\infty e_0\|^2$ and $\mu_N=\sum k|w_k|^2$.
Crucially $BW_\infty e_0=0$. For the finite unitary overlap $W_N$ and
$\rho_m=\{(m+1/4)/N\}$, the source identity is
\[
 \mathrm{Res}_N=-\operatorname{diag}(k)/N+W_N\operatorname{diag}(\rho_m)W_N^*-PBP.
\]
Hence $w^*\mathrm{Res}_Nw\ge-\mu_N/N+(1-t_N)/(4N)-t_N$.
Pairing signed momenta and using convex reciprocal-square tail integrals yields
\[
 \mu_N\le64/675,\qquad Nt_N\le T_*:=\frac{4A}{1-(3/32)^2}.
\]
For $A_N=|e_0\rangle\langle e_0|/2-|w\rangle\langle w|$, (R4.3) at
$k=0$ and the exact identity
$2\operatorname{Re}\langle A_N,\mathrm{Res}_N\rangle
=(\mathrm{Res}_N)_{00}-2w^*\mathrm{Res}_Nw$ imply
\[
 N\,2\operatorname{Re}\langle A_N,\mathrm{Res}_N\rangle
 \le A\pi+\tfrac A{16}(2+3D/2)-\tfrac12+\tfrac{128}{675}
                +2(1+1/64)T_*<0.282637.                   \tag{R4.5}
\]
This is not an absolute-value bound: the negative finite cross terms
explicitly falsify the absolute $0.283/N$ mutant. Cell contraction and C1 give
\[
 \|S_N+H_N\|_{\rm HS}^2
 \le\tfrac54-\tfrac8{\pi^2}
       +(\sqrt{0.1525}+0.466902/4)^2+0.282637/16<0.714386. \tag{R4.6}
\]

\paragraph{Validated Fourier core and the missing finite-grid aliases.}
Use the sigma-$x$ basis
$H(p)=\left(\begin{smallmatrix}\sin p\,I_8&B_p\\B_p^T&-\sin p\,I_8\end{smallmatrix}\right)$,
$B_p=(1-\cos p)I_8-S$, $S_{y,y+1}=1$. The jump is
$dP=E_{8,8}-E_{7,7}$ (zero-based indices), and
$F(p)=P(p)-(1/2-p/(2\pi))dP$ for $0<p<2\pi$.
At the 512 midpoint nodes, 256-bit Acb validated eigensystems strictly
separate eight occupied and eight unoccupied eigenvalues. The outward
matrix DFT $A_m$ obeys the exact midpoint alias formula
$A_m=\sum_{\ell\in\mathbb Z}(-1)^\ell C_{m+512\ell}$.
The v1022 bound $\|C_m\|_{\rm HS}\le K_3/|m|^3$,
$K_3=2.517464500835$, controls the entire alias error by
$K_3\zeta(3)[(512-m)^{-3}+512^{-3}]$. The certified rational caps
for $m=1,\ldots,16$ are, in order,
\begin{gather*}
 .872617451262, .119561574970, .036384250090, .013594203794,\\
 .005638853577, .002552800879, .001281003972, .000727535964,\\
 .000465728162, .000325832279, .000241157627, .000184924134,\\
 .000145312983, .000116348415, .000094610814, .000077965240.
\end{gather*}
Since $\sum_{m>16}m^{-5}\le1/(4\cdot16^4)$, these give
$2\sqrt{\sum_{m\ge1}m\|C_m\|_{\rm HS}^2}<1.783260551<1.783261$.
This is a bound on a dominant Fourier matrix, not yet on $R_N^{\rm sm}$.
Absolute summability gives the exact four-cover expression
\[
 A_N(d)=\sum_{q\notin4\mathbb Z}(1-i^q)C_{d+qN},\qquad
 \|R_N^{\rm sm}\|_{\rm HS}^2
 =\sum_{d=-(N-1)}^{N-1}(N-|d|)\|A_N(d)\|_{\rm HS}^2.
\]
Choose $q=1$ for negative $d=-N+m$ and $q=-1$ for positive $d$;
at $d=0$ choose $q=1$ once. The resulting dominant norm squared is
$4\sum_{m=1}^{N-1}m\|C_m\|^2+2N\|C_N\|^2$, bounded by the infinite core.
The other aliases have envelope
\[
 R(t)=\sum_{q\notin4\mathbb Z,\ q\ne1}
           \frac{|1-i^q|}{|q-1+t|^3},\qquad 0\le t\le1.
\]
This nonsingular series is convex. Sum $|q|\le100$ and add the full
tail $4/100^3+2/100^2$; its outward bounds at $0,1/2,1$ are denoted
$a,b,c$ (approximately $2.528674155,1.333485721,2.086708706$).
Use their actual Arb enclosures, not those abbreviated printed decimals.
The piecewise affine interpolant $L(t)$ majorizes $R(t)$.
For $f(t)=tL(t)^2$ the composite trapezoid rule on each half interval
is exact for the cubic, including the derivative jump. For every even $N$,
\[
 \begin{split}
 T_N&=\frac2N\left(\sum_{m=1}^{N-1}f(m/N)+\frac{f(1)}2\right)
       =I+D_* /(6N^2),\\
 I&=a^2/24+ab/12+b^2/3+bc/4+7c^2/24,\\
 D_*&=4b^2-2ab-a^2+5c^2-6bc.
 \end{split}
\]
Thus the alias norm is at most
$K_3\sqrt{I+|D_*|/(6\cdot16^2)}/16^2
\le0.017332234304270<0.017333$ and
$\|R_N^{\rm sm}\|_{\rm HS}<1.800594$.
Finally $R_N=R_N^{\rm sm}-(S_N+H_N)\otimes Q$, with
$\|Q\|_{\rm HS}=\sqrt2$ from two orthogonal rank-one projectors.
Using the unrounded alias enclosure with (R4.6) gives
\[
 \|R_N\|_{\rm HS}
 <1.783261+0.017332234304270+\sqrt{2(0.714386)}
 <2.995906<3.
\]
The fixed-width full collar still has two gapless edges. One-edge
microscopic identification, FE-GEN, ALG-EXH and charged MMST do not follow
from this norm theorem.

\paragraph{T3/T4: canonical DET excitations and the signed wall.}
\texttt{CHIRAL4D.MIRROR.SIGNED.CAR.01}.
\veri{v1027\_\allowbreak signed\_\allowbreak det\_\allowbreak car\_\allowbreak wall.py}
checks the following algebra and its finite countermodels.
The old Hubbard operators $b_i=|\Omega\rangle\langle e_i|$ satisfy
$\{b_i,b_j^*\}=\delta_{ij}P_\Omega+|e_j\rangle\langle e_i|$, not CAR,
and $b_i^*b_j^*=0$. For even $n\ge4$ (in particular $n=32$), write
$F=|{\rm full}\rangle\langle0|$ and define
\[
 U=I+(2^{-1/2}-1)(P_0+P_F)+(F-F^*)/\sqrt2,\qquad a_i=Uc_iU^*.
\]
This is a two-dimensional rotation and identity elsewhere. It is even;
under the inherited full-determinant singlet premise it is onsite gauge
invariant. It maps $|0\rangle$ to $\Omega=(|0\rangle+|{\rm full}\rangle)/\sqrt2$,
so the $a_i$ satisfy CAR on the entire $2^n$-dimensional Fock space and
$a_i^*\Omega=e_i$. Products of even onsite rotations preserve cross-site CAR.
The actual new parent is additive:
\[
 N_a=UNU^*=N+\tfrac n2(P_0-P_F-F-F^*),\qquad
 Q\le N_a\le nQ,\quad Q=I-P_\Omega.
\]
Thus $\Delta N_a$ has the same unique ground state and gap $\Delta$ as
$\Delta Q$, but is not the old flat parent. The high-degree local rotation
is not asserted to be a linear Bogoliubov map or an efficient Fock simulator.

Let $A=A^*$ be a fixed c-number, gauge-covariant quasilocal one-particle
operator, $\|A\|\le L$. Define
\[
 H=p^*Ap+\Delta a^*a+\lambda(a+gAp)^*(a+gAp),\quad
 M=\Delta+\lambda,\quad\beta=\lambda g^2,\quad\eta=\lambda g,
 \quad c=\lambda g^2\Delta/M.
\]
If $cL<1$ and $B=L+\beta L^2<M$, the spectral block
\[
 h(s)=\begin{pmatrix}s+\beta s^2&\eta s\\\eta s&M\end{pmatrix},\quad
 \det h(s)=Ms(1+cs)
\]
has precisely the original kernel, no extra zero, high band $\mu_+\ge M$,
and $|\mu_-|\le B$. Expansion gives
$\mu_-(s)=s+cs^2-\eta^2s^3/M^2+O(s^4)$, hence $z=1$ when $A$ is linear.
For $(\Delta,\lambda,g,L)=(3,1,1/2,2)$, $M=4$, $B=3$, $cL=3/8$.
Filling the negative band and subtracting its energy gives a nonnegative
full-Fock Hamiltonian for each fixed classical background;
$E_0\ge-Lr|\Lambda|$ with $r$ physical modes per site.
The bare mirror operator nevertheless has low-band weight of order $s^2$.
There is no hard bare-mirror response gap, nor an isolated entire
many-body mirror continuum. Operator-valued quantum-link diagonalization
does not automatically lift to a canonical Fock transformation (the
explicit SWAP mutant fails CAR).

For a free illustration only, $\gamma_i=\tau_1\otimes\sigma_i$,
$\gamma_5=\tau_3\otimes I$, $w=\sum_i(1-\cos k_i)-1$,
$R^2=1+2\sum_{i<j}(1-\cos k_i)(1-\cos k_j)\in[1,25]$ and
$D=I+(w+i\gamma\cdot\sin k)/R$ give $A=\gamma_5D$,
$A^2=2+2w/R$ and $\|A\|\le2$, with one Dirac node at zero.
The inverse-square-root series has norm ratio $12/13$ and exponential
length $2/\log(13/12)$. This four-component regulator contains opposite
Weyl sectors: it is not the demanded two-component chiral SM construction.
The common-parent dynamical gauge, measure, anomaly/index and continuum
obligations remain open.

\paragraph{T6/T8: full Gauss counts, internal cap and selection rules.}
\texttt{PARENT.GAUSS.D4.ROUND4.01}.
\veri{v1028\_\allowbreak gauss\_\allowbreak d4\_\allowbreak cap.py}
uses the full exterior character, not a low-occupation truncation.
Let $W_{16}$ consist of $Q:(\mathbf3,\mathbf2)_1$,
$u^c:(\bar{\mathbf3},\mathbf1)_{-4}$,
$d^c:(\bar{\mathbf3},\mathbf1)_2$, $L:(\mathbf1,\mathbf2)_{-3}$,
$e^c:(\mathbf1,\mathbf1)_6$, $\nu^c:(\mathbf1,\mathbf1)_0$, with charges $6Y$.
For color weights $(1,0),(0,1),(-1,-1)$ and weak weights $\pm1$, set
\[
 P(z)=\prod_{w\in W_{16}}(1+z^w),\quad
 h=1+(u+u^{-1})v^3,\quad
 \Delta_W=(1-x/y)(1-x^2y)(1-xy^2)(1-u^2).
\]
There are twelve SM16 copies per site (192 fermion modes), a three-state
Higgs space, and seam dimension four. The singlet multiplicity of a
Weyl-invariant character is $\operatorname{CT}(f\Delta_W)$, not its
zero-weight multiplicity. All representations descend to the $\mathbb Z_6$
quotient. In characteristic five,
$P^{12}=P^2(z)P^2(z^5)$, so for $A=P^2$,
\[
 \operatorname{CT}(P^{12}z^r)
   =\sum_{a+5b+r=0}A_aA_b\pmod5.
\]
This finite sparse calculation has 20771 integer two-copy coefficients
(17151 nonzero modulo five); it never enumerates $2^{192}$ states.
The full one-site Gauss dimension is $4\operatorname{CT}(P^{12}h\Delta_W)
\equiv3\pmod5$. A neutral Majorana bijects the parity sectors, each
of which has dimension $4\pmod5$.
For the nine link irreps $0,Q,u^c,d^c,L,e^c,H,\operatorname{Adj}_3,
\operatorname{Adj}_2$, whose squared dimensions sum to 137, the exact transfer
matrix is $T_{RS}=\operatorname{CT}(P^{12}h\chi_R\bar\chi_S\Delta_W)$.
An $n$-vertex open path has dimension $4^n(T^n)_{00}\pmod5$.
For $n=1,\ldots,8$ this is $3,4,3,0,0,1,4,0$.
The full open $2\times2\times2$ cube within this specified nine-irrep link
truncation (eight sites, twelve links) has
dimension $0\pmod5$. Two pairwise contraction orders give reduced integer
48544377696195; this is not the Hilbert dimension. The bound
$9^{12}4^8<2^{63}$ precludes contraction overflow.
Independent exact raising-operator ranks check 56 singlets without Higgs
($556-500$) and 120 with Higgs ($1376-1256$) for one copy; the two-copy
character result with Higgs is 561848. Thus the universal projected
non-divisibility obstruction is false. An abstract Weyl pair exists when
five divides the dimension, but locality, covariance, canonical choice,
energy compatibility and compatibility between volumes are additional demands.

On the selected even 12-dimensional pair space $|k,g\rangle$, use
$C=\operatorname{diag}(1,-i,-1,i)\otimes\operatorname{diag}(-i,-1,i)$
and $J|k,g\rangle=|-k\bmod4,4-g\rangle$. Then
\[
 \psi=(|1,3\rangle+|2,2\rangle+|3,1\rangle)/\sqrt3,\qquad
 h_{\rm sym}=I-|\psi\rangle\langle\psi|
\]
is $D_4$ invariant, has unique ground vector and gap one, Schmidt rank
three, and nonadditive HS norm $\sqrt{8/9}$. This is ordinary $D_4$ only
on that selected even space, not the whole projective fermion lift.
The invariant-vector space has dimension two, so symmetry does not
select this ray. Extending the cap by zero outside this subspace leaves
the old vacuum and does not select the full-parent state. The unitary
$X=P_\psi+e^{2\pi i/5}(I-P_\psi)$ has order five but only two eigenvalues;
it is not a Weyl clock, and its chosen order is not derived.

For the actual one-fermion phases $u_g=e^{-i\pi g/4}$, invariance of a
Majorana bilinear requires $(u_gu_h-1)M_{gh}=0$. Since $2\le g+h\le6$,
$M=0$ for this symmetric candidate. A Hermitian pair kernel instead has
commutant $\operatorname{diag}(a,b,a)$; the two transformation laws differ.
The specified squared seam-pair coupling conserves
$Q=N_\nu+2Q_{\rm seam}$ and cannot generate anomalous pairing in its
unique $Q=0$ vacuum. These symmetries have not been established for every
term of an eventual full parent. No physical neutrino texture is obtained.

\paragraph{T7: local constraints and zero-mode removal.}
\texttt{GRAV.SPIN2.LOCAL.CONSTRAINT.01}.
\veri{v1029\_\allowbreak local\_\allowbreak tensor\_\allowbreak constraints.py}
implements a quadratic target, not a TFPT embedding. For six orthonormal
symmetric tensor coordinates, let $t=(1,1,1,0,0,0)^T$, $B(k)q=q_{ij}k_j$,
$v=B^Tk$, $s=v-|k|^2t$. Set
\[
 K_p=I-tt^T/2,\quad
 K_q=|k|^2(I-tt^T)-2B^TB+tv^T+vt^T,\quad
 H=\tfrac12p^TK_pp+\tfrac12q^TK_qq,
\]
with constraints $C_0=s^Tq=0$, $C=Bp=0$. Direct polynomial algebra gives
\[
 BB^T=(|k|^2I+kk^T)/2,\quad Bs=0,\quad K_qB^T=0,\quad
 K_ps=B^Tk,\quad s^Ts=2|k|^4.
\]
The four constraints commute and $\dot C_0=k^TC$, $\dot C=0$.
For $k\ne0$ they are independent first-class constraints, leaving
$12-2\cdot4=4$ phase dimensions. Rotation to $k\parallel e_3$ and
constraint/gauge reduction leaves the two normalized TT tensors, on which
$H=\frac12\sum_{a=+,\times}(p_a^2+|k|^2q_a^2)$.
Local constraints therefore avoid putting the nonlocal TT projector into
the Hamiltonian. This is the standard linearized-gravity mechanism, not
a novel derivation of general relativity (see
\href{https://arxiv.org/abs/0710.1430}{Green--Kiriushcheva--Kuzmin}).

A finite-stencil regulator puts diagonal components on sites, off-diagonal
$ij$ components at offsets $(e_i+e_j)/2$, and vector parameters on links.
Use $\kappa_i=2\sin(ak_i/2)/a$ in every formula. Under a Brillouin shift,
the sign flip $U_i$ on vectors induces $D_i$ on symmetric tensors, and
$B(U_i\kappa)=U_iB(\kappa)D_i$, $s(U_i\kappa)=D_is(\kappa)$,
$K_q(U_i\kappa)=D_iK_q(\kappa)D_i$, $K_p=D_iK_pD_i$.
Thus the half-grid symbol is consistently glued; its only zero is $k=0$,
unlike the eight-corner $\sin k_i$ mutant.
At $k=0$, however, all constraints vanish and $p_0=rt$ gives
$H_0=-3r^2/4$: the unmodified periodic model is unbounded below.
A positive radiative target requires removal of the entire global
canonical block $(q_0,p_0)$, an extra global zero-mean prescription,
not a local constraint or a TFPT state-selection theorem.
Furthermore finite-dimensional exact CCR are impossible by the trace:
$[q_d,p_d]=i(I-dP_{d-1})$ for the explicit truncation. Collective low-energy
emergence is not excluded, but needs a controlled microscopic sequence.
No critical embedding, nonlinear constraints, universal coupling or global
gravitational state is established. The distinction between a tensor
target and controlled emergence also matters in
\href{https://arxiv.org/abs/0907.1203}{Gu--Wen}; neither cited paper fills
the TFPT embedding gap.

\paragraph{ALG2: sector generation and conditional uniform section.}
\texttt{SEAM.MMST.ALG2.FRAME.01}.
\veri{v1030\_\allowbreak alg2\_\allowbreak sector\_\allowbreak frame.py}
does not restore the refuted Vandermonde implication. Suppose, explicitly,
$H=\bigoplus_rH_r$ and the actual neutral algebra is
$\mathcal A_0=\bigoplus_r\operatorname{End}H_r$. Join $r,s$ if
$P_sTP_r\ne0$. Then $W^*(\mathcal A_0,T)=\operatorname{End}H$ precisely
when this undirected sector graph is connected: the neutral commutant is
$\bigoplus_rc_rI_r$, and each nonzero corner forces $c_r=c_s$.
The finite bicommutant theorem proves the result. Genuine superselection
components must be treated separately, not bridged by forbidden odd fields.

For a fixed $D$-dimensional low-energy space choose actual local words
$Y_{j,N}$, $j=1,\ldots,D^2$, with $\|Y_{j,N}\|\le M$.
Let $\Phi_N$ have columns $\operatorname{vec}(P_NY_{j,N}P_N)$.
If $\sigma_{\min}(\Phi_\infty)=\sigma>0$ and
$\|\Phi_N-\Phi_\infty\|_2\le\sigma/2$, then solving
$\Phi_Na=\operatorname{vec}X$ gives, for $\|X\|\le1$,
\[
 P_Ny_NP_N=X,\quad y_N=\sum_ja_jY_{j,N},\quad
 \|a\|_2\le2\sqrt D/\sigma,\quad
 \|y_N\|\le2MD^{3/2}/\sigma.
\]
This supplies a uniform section only under the stated microscopic word,
norm and noncollapsing-frame hypotheses. Compression does not preserve
products ($PABP$ need not equal $PAPBP$), and double-cutoff convergence
does not control $|N\rangle\langle0|$ on a fixed input cutoff.
For ALG2 one needs a joint real-linear frame for
$L_N(y)=(P_FyP_E,P_Fy^*P_E)$ and uniform high-output tails of both
$x_N,y_N$ and their adjoints. If the joint error and each pair of tails
are at most $\varepsilon$, the desired one-sided sum is at most
$3\varepsilon$. Neutral generation, charged corners, actual local words,
their scaling limits, and these tails are not yet proved for the one-edge
TFPT microscopic algebra. They, together with the common-parent gates,
are the precise remaining research obligations.
